% introduction.tex
\section{Introduction}\label{intro}
The integrity and efficiency of public procurement processes are crucial for ensuring the effective use of public funds, supporting economic development, and maintaining public trust \citep{Bosio2022}. Public procurement can account for a significant share of GDP (around 12\% globally and up to 20\% in developing countries), making its proper functioning vital for national economies \citep{Bosio2022}. However, anti-competitive practices in procurement --- such as collusion among bidders or corruption and favoritism by officials --- can dramatically increase costs for governments and taxpayers, reduce the quality of public services, and undermine fair competition \citep{Sokol2017}. These practices stifle innovation and disadvantage honest businesses, ultimately eroding citizens' confidence in government.

Detecting collusion or other anti-competitive behavior in procurement is inherently challenging due to its covert nature and the complexity of bidding processes. Collusive arrangements are often hidden behind layers of corporate secrecy and subtle bidding patterns, making them difficult to uncover without advanced analytical tools \citep{MarshallMarx2012}. While a substantial body of research on procurement in developing countries focuses on corruption and transparency issues, much less attention has been devoted to identifying collusive behavior and bid rigging. Notable exceptions include recent studies that attempt to detect collusion in auction markets \citep[e.g.][]{kawai2022study, baranek2021detection}. This gap is particularly relevant in contexts like Brazil, where large-scale procurement programs operate. In the state of São Paulo, the economic hub of Brazil, billions of dollars are tendered annually through electronic platforms. Recognizing the risks, Brazil has introduced reforms (such as stricter procurement regulations in the 2010s) to curb corruption and collusion \citep{OECD2016}. Nevertheless, developing reliable methods to reveal anti-competitive patterns remains a pressing need.

In this paper, we address this need by applying the regression discontinuity design (RDD) methodology proposed by \cites{kawai2023using} to a comprehensive dataset of 3.8~million public procurement auctions in São Paulo state from 2009 to 2019. This approach focuses on auctions where the winning bid barely outbids the runner-up, allowing us to examine whether the outcomes around these narrow thresholds exhibit signs of collusion or favoritism. In particular, we investigate whether firms that win by a very small margin show unusual advantages --- such as a higher likelihood of having won their previous auction (incumbency) or a larger backlog of recent contracts --- compared to the firms they narrowly beat. If such patterns are present, they would signal that factors other than competitive pricing are influencing the allocation of contracts.

Our analysis reveals clear signs of anti-competitive outcomes. Firms that barely win an auction tend to have markedly stronger records of recent success than those that barely lose. In other words, a firm just winning by a narrow margin is significantly more likely to have been the incumbent (winner of its last participated auction) than a firm that just loses. Moreover, narrowly losing firms exhibit substantially smaller outstanding workloads (backlogs) in the months prior to the auction compared to narrowly winning firms. These findings suggest that certain firms accumulate consecutive wins and more ongoing contracts, consistent with preferential treatment or insider advantages.

By highlighting these patterns, our paper contributes to the literature on procurement and competition. We provide evidence from a large developing-country market that anti-competitive dynamics --- potentially driven by political favoritism or other forms of undue influence --- can skew procurement outcomes. Unlike classical cartel behavior where winning opportunities are rotated among firms, the pattern we document points toward a scenario in which a particular firm or set of firms is repeatedly favored. Our results underscore the importance of monitoring not only overt corruption but also subtle forms of favoritism and collusion in public procurement. The remainder of the paper is structured as follows. Section~\ref{coll} lays out our empirical strategy for detecting anti-competitive behavior. Section~\ref{hist} provides the institutional background of São Paulo's public procurement system, and Section~\ref{data} describes the data and key variables. In Section~\ref{res}, we present the main empirical results, with additional analyses (including heterogeneous effects and robustness checks) in Section~\ref{res_ar}. Finally, Section~\ref{concl} concludes with a discussion of the implications of our findings.
